% !TeX root = ../navier-stokes-manuscript.tex
% ============================================================
% 1.1 The Regularity Problem
% ============================================================
\subsection{The Regularity Problem}\label{sub:TheRegularityProblem}

The three-dimensional Navier--Stokes regularity problem asks whether a smooth solution can lose its smoothness in finite time.
Fix smooth, divergence-free, finite-energy initial data \(u_0\), fix \(\nu>0\), and let \((u,p)\) be the corresponding maximal classical solution on \([0,T_*)\).
We want to prove that \(T_*\) cannot be finite.

This paper treats the whole-space problem on \(\mathbb R^3\), with zero external force and rapidly decaying initial data.
The velocity \(u(x,t)\in\mathbb R^3\) and pressure \(p(x,t)\in\mathbb R\) satisfy
\[
  \partial_tu+(u\cdot\nabla)u+\nabla p
  =
  \nu\Delta u,
  \qquad
  \nabla\cdot u=0,
  \qquad
  u(x,0)=u_0(x).
\]
Classical local theory gives a unique smooth solution for a short time.
The number \(T_*\in(0,\infty]\) is the largest time up to which that solution can remain classical, so the whole-space existence-and-smoothness problem is exactly
\[
  T_*=\infty
\]
for every admissible \(u_0\) \parencite{Fefferman2000}.

\divider{Smoothness}

\begin{definition}[Smoothness]\label{def:Smoothness}
A velocity--pressure pair \((u,p)\) is smooth through a finite time \(T>0\) when
\[
  u\in C^\infty\!\left(\mathbb R^3\times[0,T];\mathbb R^3\right),
  \qquad
  p\in C^\infty\!\left(\mathbb R^3\times[0,T]\right).
\]
Equivalently, every mixed spatial and temporal derivative of every finite order exists and is continuous through \(t=T\).
The pair is smooth for all time when it is smooth through every finite \(T>0\).
\end{definition}

The word smooth has a precise all-orders meaning.
It does not require one numerical constant to bound every derivative order simultaneously.
It requires each finite derivative order to remain well defined and continuous whenever that order is examined.

\divider{Regularity}

\begin{definition}[Regularity]\label{def:Regularity}
Regularity describes the type and degree of control available for a solution.
In the classical scale, \(C^0\) means continuity, \(C^1\) adds continuous first derivatives, and \(C^k\) requires all mixed derivatives through total order \(k\).
Smoothness is the all-orders case \(C^\infty\).
Sobolev and mixed space--time bounds are other forms of regularity, expressed through norms that measure integral control of the field and its derivatives.
\end{definition}

A regularity statement must name its level.
Continuity of the velocity does not by itself control its gradient, and control of finitely many derivatives does not by itself give \(C^\infty\).
The historical continuation criteria below matter because they identify levels of regularity strong enough to recover all remaining finite orders from the equations.

\divider{Continuation and Maximal Time}

\begin{definition}[Continuation]\label{def:Continuation}
Suppose \((u,p)\) solves the Navier--Stokes equations on \(\mathbb R^3\times[0,T)\).
A continuation past \(T\) is a velocity--pressure pair \((\widetilde u,\widetilde p)\) solving the same equations with the same viscosity on \(\mathbb R^3\times[0,T+\varepsilon)\) for some \(\varepsilon>0\), agreeing with \((u,p)\) before \(T\), and belonging to the required regularity class on the larger interval.
\end{definition}

For the Clay problem, the required class is \(C^\infty\).
The adjective maximal means that the smooth solution has no such continuation beyond \(T_*\).
Thus \(T_*<\infty\) can occur only if the solution loses enough regularity near \(T_*\) to defeat every applicable local restart.
